\section{Conclusão}
\label{sec:conclusao}

Neste artigo, estudamos os leilões sob a perspectiva dos jogos bayesianos, mostrando como a Transformação de Harsanyi permite modelar e resolver disputas estratégicas nas quais cada comprador possui informação privada sobre sua disposição a pagar. Partindo dos conceitos de tipo privado, crença a priori comum (\textit{prior} comum), estratégias funcionais e Equilíbrio de Nash Bayesiano, distinguimos os leilões de Valores Privados Independentes dos leilões de Valor Comum, nos quais a vitória sem o devido desconto estatístico expõe o licitante mais otimista à Maldição do Vencedor.

No modelo analítico resolvido para dois licitantes simétricos com valorações uniformes em $[0, 1]$, demonstramos que o Leilão de Primeiro Preço induz um sombreamento de lance (\textit{bid shading}) de $50\%$ ($b^*(v_i) = v_i/2$), enquanto o Leilão de Segundo Preço de Vickrey torna a revelação sincera ($b^*(v_i) = v_i$) uma estratégia fracamente dominante. Comprovamos analiticamente e por meio de uma simulação de Monte Carlo de $10.000$ leilões em Python que ambos os formatos convergem para a mesma receita esperada teórica $\mathbb{E}[R] = \frac{n-1}{n+1} = \frac{1}{3}$, confirmando o Teorema da Equivalência de Receita.

Por fim, o estudo dos leilões de anúncios online (Google Ads e \textit{Ad Exchanges}) ilustrou como a teoria orienta decisões reais da indústria: a migração do primeiro para o segundo preço em 2002 eliminou a instabilidade dos ciclos em dente de serra, enquanto a transição de volta ao primeiro preço em 2019 respondeu à competição simultânea em múltiplos estágios do \textit{Header Bidding} e à busca por transparência, consolidando o uso moderno dos algoritmos de \textit{bid shading}.
